\chapter{リターンとリスク}\label{ch:ret}

この部（第 II 部）では、後の章を読むのに必要な金融工学の基礎を、このプロジェクトで実際に使う形で説明する。この章は、すべての出発点になるリターンとリスクの測り方である。

\section{価格からリターンへ}

\subsection{単純リターン}

株価を $P_t$、その日の配当を $D_t$ とすると、1 日の\term{単純リターン}は
\begin{equation}
R_t = \frac{P_t + D_t}{P_{t-1}} - 1
\end{equation}
である。配当を含めたものを\term{トータルリターン}と呼ぶ。株式分割があると $P_t$ が不連続に変わるので、実務では分割と配当を過去にさかのぼって調整した\term{調整後終値} $\tilde P_t$ を使い、$R_t = \tilde P_t / \tilde P_{t-1} - 1$ とする。

\begin{defbox}[定義：ポートフォリオのリターン]
比率 $w_i$（合計 1）で銘柄を持つポートフォリオの 1 期間のリターンは、銘柄のリターンの加重平均
\[ R_{p,t} = \sum_i w_i R_{i,t} = w\T R_t \]
である。単純リターンは銘柄の間で足せる。この性質のため、ポートフォリオの計算には単純リターンを使う。
\end{defbox}

\subsection{対数リターンと期間のつなぎ方}

\term{対数リターン} $r_t = \log(1+R_t)$ は、期間の間で足せる。$T$ 日の累積のリターンは
\[ 1 + R_{0\to T} = \prod_{t=1}^{T}(1+R_t), \qquad r_{0\to T} = \sum_{t=1}^{T} r_t \]
である。1 か月はおよそ 21 取引日、1 年は 252 取引日として数える。

\subsection{期間の単位を変える}

日次のリターンの平均を $\mu_d$、標準偏差を $\sigma_d$ とし、日々のリターンが独立で同じ分布に従うと仮定すると、$n$ 日の
\begin{equation}
\mu_n \approx n\,\mu_d, \qquad \sigma_n \approx \sqrt{n}\,\sigma_d
\label{eq:sqrt-time}
\end{equation}
になる。平均は日数に比例し、標準偏差は日数の平方根に比例する。

\begin{exbox}
日次の標準偏差が 2\% の銘柄は、1 か月（21 日）で $2\% \times \sqrt{21} \approx 9.2\%$、1 年で $2\% \times \sqrt{252} \approx 31.7\%$ になる。日次の平均が 0.05\% なら、1 か月で約 1.05\%、1 年で約 12.6\% である。
\end{exbox}

\begin{notebox}[単位をそろえる]
平均と標準偏差、共分散、リスク回避度、無リスク金利は、すべて同じ期間の単位でなければならない。単位が混ざると、結果が桁で狂う。このプロジェクトでは実際に 2 回これを踏んだ。
\begin{itemize}
\item 元論文の実装は、日次のリターンから年率 2\% の無リスク金利をそのまま引いていた。均衡リターンが「毎日 1〜2\% 下がる」値になった（総合報告書 1）。
\item 実験 7・9 は、1 か月持つポートフォリオの見通しを日次の $\sigma$ で作り、期間に合った大きさの $\sqrt{21}\approx4.6$ 倍の $\alpha$ を入れていた（X05）。TE の上限が効いていたので、結果は変わらなかった。
\end{itemize}
\end{notebox}

\section{超過リターン}

リターンそのものより、何かと比べた差のほうが意味を持つことが多い。比べる相手によって名前が違う。

\begin{longtable}{P{3.4cm}P{5.0cm}P{6.0cm}}
\toprule
\textbf{名前} & \textbf{式} & \textbf{使う場面} \\
\midrule
\endhead
\bottomrule
\endfoot
無リスク金利に対する超過リターン & $R_i - R_f$ & CAPM、シャープ・レシオ。$R_f$ は 3 か月物の米国債の利回り（\code{\textasciicircum IRX}） \\
アクティブ・リターン & $R_p - R_b$ & ベンチマーク（\SP{} や時価総額加重）と比べた運用の成績。TE と情報レシオの元になる（\chapref{ch:active}） \\
異常リターン（アブノーマル・リターン） & $R_i - \E[R_i \mid \text{市場}]$ & イベントスタディ。このプロジェクトでは、同じセクターの ETF のリターンを引いた差（\chapref{ch:event}） \\
\end{longtable}

\section{リスクの測り方}

\subsection{分散と標準偏差}

リターンのばらつきを\term{リスク}と呼ぶ。最も基本的な測り方は分散 $\sigma^2 = \E[(R-\mu)^2]$ と、その平方根の標準偏差 $\sigma$（\term{ボラティリティ}）である。過去 $T$ 日のデータから
\[ \hat\sigma^2 = \frac{1}{T-1}\sum_{t=1}^{T}(R_t - \bar R)^2 \]
で推定したものを\term{実現ボラティリティ}と呼ぶ。

\subsection{予想されたボラティリティ：VIX}

\term{VIX} は、\SP{} のオプションの価格から逆算した、市場が予想する翌 30 日の \SP{} のボラティリティ（年率、\%）である。過去のデータから計算する実現ボラティリティと違い、市場の\textbf{予想}を表す。ふだんは 12〜20 程度で、2020 年 3 月のコロナの急落では 80 を超えた。VIX は実現ボラティリティより平均して少し高い（オプションの売り手が受け取る保険料の分）。

\subsection{下側のリスク：VaR と CVaR}

分散は上振れと下振れを同じに扱う。下振れだけを見る測り方に、\term{VaR}（バリュー・アット・リスク）と \term{CVaR}（条件付き VaR、期待ショートフォール）がある。損失 $L=-R$ の分布で、信頼水準 $\alpha$（たとえば 95\%）に対して
\[ \VaR_\alpha = \inf\{\ell : \Pr(L \le \ell) \ge \alpha\}, \qquad \CVaR_\alpha = \E[L \mid L \ge \VaR_\alpha] \]
である。VaR は「$1-\alpha$ の確率で起きる損失の境目」、CVaR は「その境目を超えたときの損失の平均」。CVaR は凸で、最適化に組み込みやすい（第 IV 部「CVaR で最適化する」）。

\section{分散投資と共分散}

2 つの銘柄のリターンの\term{共分散}は $\Cov(R_i,R_j) = \E[(R_i-\mu_i)(R_j-\mu_j)]$、\term{相関}は $\rho_{ij} = \Cov(R_i,R_j)/(\sigma_i\sigma_j)$ である。$N$ 銘柄の共分散をまとめた $N\times N$ の行列を $\Sigma$ と書くと、ポートフォリオの分散は
\begin{equation}
\sigma_p^2 = w\T \Sigma w = \sum_i\sum_j w_i w_j \Sigma_{ij}
\label{eq:port-var}
\end{equation}
になる。相関が 1 より小さければ、ポートフォリオのリスクは個別のリスクの加重平均より小さくなる。これが\term{分散投資}の効果である。

\begin{exbox}
標準偏差 30\% の銘柄を $N$ 個、同じ比率で持ち、どの 2 銘柄の相関も 0.3 だとする。ポートフォリオの分散は $\sigma^2[\rho + (1-\rho)/N]$ なので、$N=500$ では $0.09\times(0.3+0.7/500)\approx 0.0271$、標準偏差は約 16.5\% になる。いくら銘柄を増やしても、相関の分 $\sigma\sqrt{\rho}\approx16.4\%$ は消えない。これが市場全体のリスク（\term{システマティック・リスク}）である。
\end{exbox}

共分散の行列を\textbf{どう推定するか}は、それ自体が大きな問題で、次の章で扱う。

\section{ベータと CAPM}

\subsection{ベータ}

銘柄 $i$ の\term{ベータ}は、市場のリターン $R_m$ に対する感応度
\begin{equation}
\beta_i = \frac{\Cov(R_i, R_m)}{\Var(R_m)}
\end{equation}
である。市場が 1\% 動くと、その銘柄は平均して $\beta_i$\% 動く。ベータが 1 より大きい銘柄（半導体など）は市場より大きく動き、小さい銘柄（公益、生活必需品など）は小さく動く。過去 252 日の日次リターンを \SP{} に回帰した傾きで推定することが多い。

\subsection{CAPM}

\term{CAPM}（資本資産価格モデル）は、すべての投資家が平均分散で考え、同じ予想を持つときに成り立つ均衡の式
\begin{equation}
\E[R_i] - R_f = \beta_i\,(\E[R_m] - R_f)
\label{eq:capm}
\end{equation}
である。$\E[R_m]-R_f$ を\term{市場のリスクプレミアム}と呼ぶ。CAPM が言っているのは、「報われるのは市場全体のリスクだけで、分散で消せる個別のリスクには上乗せが付かない」ということである。

Black--Litterman モデルの出発点である\term{均衡リターン} $\Pi = \delta\Sigma\wmkt$ は、CAPM を逆向きに使ったものである（第 III 部）。市場全体が時価総額の比率 $\wmkt$ を持っていることを「均衡」とみなし、それが最適になるような期待リターンを逆算する。

\section{指数と時価総額加重}

\SP{} は、米国の大型株約 500 銘柄を、浮動株調整後の時価総額で加重した指数である。構成銘柄は委員会が定期的に入れ替える。時価総額加重のポートフォリオは、
\begin{itemize}
\item 株価が動いても比率が自動で追従するので、組み替えの取引がほとんど要らない（費用が小さい）。
\item 市場全体の平均そのものなので、全員がアクティブに運用しても、平均はこのリターンになる。
\end{itemize}
という性質を持つ。アクティブ運用の成績は、この時価総額加重と比べて測る（\chapref{ch:active}）。

\begin{notebox}[生存者バイアス]
「今の構成銘柄」で過去を計算すると、途中で外れた銘柄（業績が悪化して外れた、買収された、上場を廃止した）が抜け落ちる。生き残った銘柄だけで過去を見ると、成績が良く見える。これを\term{生存者バイアス}と呼ぶ。このプロジェクトでは、各判断日にその時点の構成銘柄を使うことで避けた。ただし yfinance には上場を廃止した銘柄の株価がないものがあり、その銘柄はデータから欠ける（\chapref{ch:pit}）。
\end{notebox}

\section{リスクプレミアムとファクター}

CAPM の後、市場のベータ以外にも、平均して報われるリスクがあることが報告されてきた。代表的なものを\term{ファクター}と呼ぶ。

\begin{longtable}{P{2.6cm}P{5.6cm}P{6.2cm}}
\toprule
\textbf{ファクター} & \textbf{中身} & \textbf{よく言われる理由} \\
\midrule
\endhead
\bottomrule
\endfoot
市場 & 株式 − 無リスク金利 & 株式のリスクを負うことへの報酬 \\
規模（SMB） & 小型株 − 大型株 & 流動性が低く、倒産しやすい \\
割安（HML） & 簿価時価比率の高い株 − 低い株 & 業績の悪い会社のリスク、または投資家の過剰な悲観 \\
モメンタム（UMD） & 過去 12 か月（直近 1 か月を除く）に上がった株 − 下がった株 & 投資家の反応の遅れ（行動ファイナンス） \\
低ボラティリティ & 値動きの小さい株 − 大きい株 & レバレッジを使えない投資家が高ベータを買いすぎる \\
収益性・投資 & 利益率が高い・投資が控えめな会社 & 割引率と利益の関係（Fama--French の 5 ファクター） \\
\end{longtable}

ファクターのリターンの日次のデータは、Kenneth French のサイトで公開されている。このプロジェクトでも、追加のデータとして取得した（\code{factor\_returns}）。

アクティブ運用の多くは、結局これらのファクターのどれかに傾けることで成績を上げている、という見方がある（Pedersen『Efficiently Inefficient』）。市場が「落とした」リスクプレミアムを拾うのが、アクティブ運用の利益の源泉の 1 つである。実験 G は、この考え方で LLM に「いまどのリスクが報われるか」を聞いたものだった。

\section{効率的市場仮説}

\term{効率的市場仮説}は、「価格は手に入る情報をすでに反映している」という考え方で、反映される情報の範囲で 3 つに分けられる。

\begin{description}[leftmargin=1.5em,style=nextline]
\item[弱度] 過去の価格の情報は、すでに価格に入っている。過去の値動きから将来は当てられない。
\item[準強度] 公開されたすべての情報（決算、ニュース、提出書類）は、公開と同時に価格に入る。
\item[強度] 公開されていない情報まで価格に入っている。
\end{description}

現実の市場は、準強度に近いが完全ではない。よく知られた例外が、決算発表の後に株価がしばらく同じ向きに動き続ける\term{発表後のドリフト}（PEAD）である（\chapref{ch:event}）。この研究の問い 3（LLM の読み取りで儲かるか）は、「決算発表文の文章の情報に関して、市場は準強度の効率性を満たしているか」という問いと同じ形をしている。結果は、少なくとも日次の終値で見る限り「満たしている」だった。

\begin{projbox}
\begin{itemize}
\item \textbf{リターン}：DB の株価は yfinance の調整後終値。日次の単純リターンから、月次・期間のリターンを積で作る（第 VI 部）。
\item \textbf{LLM への入力}：元論文と実行 1 は、判断月の日次リターン（約 20 個）をそのまま LLM に渡し、翌月の平均日次リターンを答えさせた（第 VII 部）。
\item \textbf{超過リターン}：実験 2 以降の目的変数は、銘柄のリターンからセクター ETF のリターンを引いた 2 取引日の異常リターン（\chapref{ch:event}）。
\item \textbf{単位}：均衡リターンの無リスク金利の誤り（総合報告書 1）と、$\alpha$ の $\sqrt{21}$ の誤り（X05）。
\item \textbf{VIX}：実験 G・H05 の入力、X04 の共分散の大きさの合わせ方。
\item \textbf{ファクター}：実験 G のリスクの種類、追加データの Fama--French。
\item \textbf{時価総額加重}：BLM の均衡の比率で、実験 7〜9 のベンチマーク。
\end{itemize}
\end{projbox}
